\title{Teriyakinudler med hakket gris og peanuts}
\PrepTime{20}
\CookingTime{10}
\CookingTempe{}
\TypeCooking{Hurtige måltider\index[nemlig]{Hurtig!Teriyakinudler med hakket gris og peanuts}}
\NbPerson{4}
\Image{0 100 800 900}{images/nemlig/teriyakinudler-m-hakket-gris-peanuts.png}

\begin{ingredient}
    \begin{main}
        \item 1 pk. Nudler\index[ingredient]{Nudler}
        \item Salt\index[ingredient]{Salt}
        \item 40g Ingefær\index[ingredient]{Ingefær}
        \item 80g Peanuts\index[ingredient]{Peanuts}
        \item Olie\index[ingredient]{Olie}
        \item 1 pk. Hakket Svinekød\index[ingredient]{Hakket Svinekød}
        \item Peber\index[ingredient]{Peber}
        \item 1 pk. Grøntsagsblanding\index[ingredient]{Grøntsagsblanding}
        \item 8 spsk. Teriyakiglaze\index[ingredient]{Teriyakiglaze}
    \end{main}
\end{ingredient}
\begin{recipe}
    \step{Kog nudler i letsaltet vand i ca. 3--4 min., til de er møre. Si vandet fra, men gem en kop af kogevandet til at justere saucen med. Skyl dem under koldt vand og lad dem dryppe af.}
    \step{Skræl ingefæren, og riv fint på et rivejern.}
    \step{Hak peanuts fint.}
    \step{Varm neutral olie op i en wokpande eller stor stegepande. Steg ingefæren ved høj varme i ca. 1 min.}
    \step{Tilsæt det hakkede grisekød, og steg yderligere i 4--5 min. (fortsat ved høj varme), til det er gennemstegt og brunet godt af. Krydr med friskkværnet peber.}
    \step{Vend grøntsagsblandingen i wokken, og steg i 3--4 min. (fortsat ved høj varme) – rør jævnligt.}
    \step{Skru ned til middelvarme. Tilsæt de kogte nudler og hæld teriyakiglazen over. Vend wokretten godt rundt, og steg i yderligere 2--3 min., indtil alt er gennemvarmt og saucen er jævnt fordelt. Justér evt. konsistensen med lidt af det gemte kogevand, så nudlerne ikke bliver tørre.}
    \step{Servér den hurtige wok med teriyakinudler og top med hakkede peanuts (se tip!).}
\end{recipe}
\begin{notes}
    \textbf{Tip!} Top gerne retten med friske krydderurter, hvis du har nogen stående i vindueskarmen - især koriander er oplagt til denne ret.
\end{notes}

\setcounter{step}{0}
\setcounter{substep}{0}

\clearpage